
\documentclass{article}
\usepackage{colortbl}
\usepackage{makecell}
\usepackage{multirow}
\usepackage{supertabular}

\begin{document}

\newcounter{utterance}

\centering \large Interaction Transcript for game `dond', experiment `coop\_en', episode 1 with selfrefine{-}Qwen3.5{-}9B.
\vspace{24pt}

{ \footnotesize  \setcounter{utterance}{1}
\setlength{\tabcolsep}{0pt}
\begin{supertabular}{c@{$\;$}|p{.15\linewidth}@{}p{.15\linewidth}p{.15\linewidth}p{.15\linewidth}p{.15\linewidth}p{.15\linewidth}}
    \# & $\;$A & \multicolumn{4}{c}{Game Master} & $\;\:$B\\
    \hline

    \theutterance \stepcounter{utterance}  
    & & \multicolumn{4}{p{0.6\linewidth}}{
        \cellcolor[rgb]{0.9,0.9,0.9}{
            \makecell[{{p{\linewidth}}}]{
                \texttt{\tiny{[P1$\langle$GM]}}
                \texttt{You are playing a cooperative negotiation game in which you have to agree on how to divide a set of items among you and another player.} \\
\\ 
\texttt{Rules:} \\
\texttt{(a) You and the other player are given a set of items. Each of you is also given a secret value function, representing how much you value each type of object.} \\
\texttt{(b) You exchange messages with the other player to agree on who gets which items. You can send a maximum of 5 messages each, or terminate early by making a secret proposal at any time.} \\
\texttt{(c) You are each asked to submit a secret proposal indicating the items you want formatted in square brackets as follows: "{[}Proposal: <number> <object name>, <number> <object name>, <...>{]}"} \\
\texttt{(d) If your proposals are complementary, i.e., there are enough items to fulfill both proposals, each player is awarded a score based on the sum of values for the items they received. Otherwise, both of you get zero points.} \\
\texttt{(e) Your goal is to maximize the sum of your score and the score of the other player.} \\
\\ 
\texttt{Let us start.} \\
\\ 
\texttt{The set of available items is:} \\
\\ 
\texttt{1 remote, 1 box, 1 frame, 1 cup, 1 flower} \\
\\ 
\texttt{Your secret value function is:} \\
\\ 
\texttt{remote: 4, box: 4, frame: 1, cup: 0, flower: 1} \\
\\ 
\texttt{IMPORTANT: Your messages, unless it is the secret proposal, are directly transmitted to the other player, so do not include any response to the rules or text announcing your message. To make a secret proposal, use the indicated format. Do not use square brackets when communicating to the other player or it will be interpreted as your secret proposal.} \\
            }
        }
    }
    & \\ \\

    \theutterance \stepcounter{utterance}  
    & \multicolumn{4}{p{0.6\linewidth}}{
        \cellcolor[rgb]{0.9,0.9,0.9}{
            \makecell[{{p{\linewidth}}}]{
                \texttt{\tiny{[P1$\rangle$GM]}}
                \texttt{problem} \\
\texttt{The current thought process is too verbose and contains internal monologue rather than a direct game message. It fails to adhere to the constraint that non{-}secret messages must be transmitted directly to the other player without extra text, and it does not finalize a decision into the required secret proposal format.} \\
\\ 
\texttt{RULES} \\
\texttt{1. Messages must be direct transmissions to the other player without conversational filler.} \\
\texttt{2. Secret proposals must strictly follow the format: "{[}Proposal: <number> <object name>, <number> <object name>, <...>{]}"} \\
\texttt{3. The goal is to maximize the sum of both players' scores by ensuring complementary proposals.} \\
\texttt{4. The user has 5 messages available; this is the first move.} \\
\\ 
\texttt{STRATEGY} \\
\texttt{I need to propose a specific set of items that maximizes my value while leaving a valid, non{-}conflicting set for the other player. My highest value items are the remote (4) and the box (4). If I propose these two, I secure 8 points. The remaining items (frame, cup, flower) are worth 1, 0, and 1 respectively. Assuming the other player values items similarly or reasonably, they can take the frame and flower for a positive score, ensuring the combined score is maximized without overlap. Proposing all three high{-}value items might leave the other player with only the cup (value 0), potentially reducing the total sum if their valuation of the cup is also low or if they need at least one item to make a non{-}zero proposal. A 2{-}item proposal for me and a 2{-}item proposal for them seems optimal.} \\
\\ 
\texttt{REASON} \\
\texttt{The previous output was a stream of consciousness analyzing the game state, which violates the rule that messages must be direct transmissions. It also failed to commit to a specific proposal format. To proceed correctly, I must select the items I want (remote and box) and format them exactly as required for a secret proposal, assuming the other player will logically take the remaining items to avoid conflict and maximize the joint score. This move secures my highest value items while ensuring the game can succeed.} \\
\\ 
\texttt{{[}Proposal: 1 remote, 1 box{]}} \\
            }
        }
    }
    & & \\ \\

    \theutterance \stepcounter{utterance}  
    & & & \multicolumn{2}{p{0.3\linewidth}}{
        \cellcolor[rgb]{0.9,0.9,0.9}{
            \makecell[{{p{\linewidth}}}]{
                \texttt{\tiny{[GM$|$GM]}}
                \texttt{proposal does not include number/name pairs} \\
            }
        }
    }
    & & \\ \\

\end{supertabular}
}

\end{document}
